% ===================================================================
%  LENTELĖ Nr. 1 — Mokykla. — Licėjus. — Klasė.
%
%  Ceci n'est PAS une transcription : c'est une traduction, faite en
%  2026, en lituanien standard. D'ou trois differences avec texto/io
%  et texto/fr, et trois seulement :
%
%   * pas de \nl ni de \cc ;
%   * pas de \begin{VUpage} ;
%   * le gras se pose sur le TERME seul, ponctuation dehors.
%
%  Les cles « %%K » sont celles de texto/io, macro pour macro pour
%  l'apparat, et les renvois \textsuperscript{(n)} visent les MEMES
%  objets de la planche.
%
%  QUATORZIEME COLONNE DES DIX-SEPT, et la premiere langue BALTE du
%  volume. Aucune des treize precedentes n'appartient a cette
%  famille ; le slave en est le voisin le plus proche, et c'est
%  precisement ce qui rend le tcheque et l'ukrainien de MAUVAIS
%  modeles ici — assez proches pour qu'on les recopie sans y penser,
%  assez loin pour que la copie soit fausse.
%
%  LE LITUANIEN DECLINE SEPT CAS COMME LE TCHEQUE, MAIS IL NE RANGE
%  PAS SES MOTS DANS LE MEME ORDRE, et c'est le point a tenir des
%  maintenant. Le tcheque POSTPOSE son genitif comme l'ido — « dveře
%  hradu » — de sorte que la colonne treize n'a pas paye un seul
%  tour d'inversion. Le lituanien l'ANTEPOSE : « mokyklos kiemas »,
%  la cour de l'ecole, met l'ecole devant. C'est la construction qui
%  a coute treize tours au finnois et vingt-neuf au basque, et elle
%  revient ici dans une langue indo-europeenne, ce qui ne s'etait pas
%  encore vu. Elle ne coute rien tant qu'UN SEUL des deux mots porte
%  un renvoi — le critere ne compte que les paires ou les DEUX en
%  portent un — et le tableau 1 en donne trois exemples inoffensifs :
%  « muzikos ir piešimo klases », « rankų darbų dirbtuvės », «
%  fizikos ir chemijos prietaisai ». Aucun n'appareille deux renvois.
%  On les signale tout de meme, parce que c'est la meme construction
%  qui, ailleurs, en appareillera deux.
%
%  ET LE GRAS PORTE LA FORME QUE LA PHRASE DEMANDE, non celle du
%  dictionnaire : « ant kėdės » et non « kėdė ». Meme parti qu'en
%  russe, en ukrainien, en basque et en tcheque.
%
%  L'ORDINAL A DEUX ETATS. « pirma serija » et « pirmoji serija »
%  disent tous deux la premiere serie ; la forme longue, dite
%  pronominale, designe un objet deja connu, et c'est elle que
%  l'usage met devant un nom precis. On ecrit « PIRMOJI SERIJA » a
%  l'apparat et l'on declare les deux etats a outils/html.py, comme
%  on avait declare les deux etats du premier ordinal irlandais.
%
%  ET LE DANGER DE CETTE COLONNE N'EST NI LE RUSSE NI LE POLONAIS :
%  c'est le LETTON. Les deux langues baltes s'ecrivent presque de la
%  meme facon, et ce qui les separe a l'oeil tient a une poignee de
%  signes — les macrons ā, ē, ī et les lettres ģ, ķ, ļ, ņ — que le
%  lituanien n'ecrit jamais. Attention au piege le plus fin : le
%  lituanien ecrit ė, un E SURPOINTE, la que le letton ecrit ē, un E
%  MACRON ; l'oeil ne les distingue pas, le code du caractere si.
%  En revanche le ū a macron appartient AUX DEUX, et l'ogonek — ą, ę,
%  į, ų — est lituanien de plein droit : le balayage qui cherche le
%  polonais ne doit pas le prendre pour un ę polonais.
%
%  ENFIN UNE REGLE DE BALAYAGE D'UNE ESPECE NOUVELLE. Les treize
%  colonnes precedentes cherchaient toutes des signes ETRANGERS
%  PRESENTS. Le lituanien permet l'inverse : son alphabet ne contient
%  ni Q, ni W, ni X, et un mot qui en porte une vient forcement
%  d'ailleurs. Un controle par ABSENCE, et il attrape ce qu'aucun
%  controle par presence n'attrape — « taxi » ecrit tel quel au lieu
%  de « taksi ».
%
%  ET LES NOMS DES NOTES RESTENT CEUX DE L'IDO, comme dans les treize
%  colonnes precedentes : la ligne du fac-simile donne l'equivalence
%  que l'ido pose, non le solfege du lecteur.
%
%  « postigo » (t01, renvoi 78) EST UN VASISTAS, ET LA PLANCHE LE DIT.
%  Le mot a ete verifie sur la gravure en 2026, non sur les
%  dictionnaires : la verriere de la classe a deux rangees de carreaux,
%  et dans la rangee HAUTE deux petits vantaux sont graves battants,
%  gonds sur un montant vertical, ouverts vers l'interieur. C'est le
%  seul ouvrant de la verriere, et il ne sert qu'a aerer.
%
%  D'OU CE QU'IL N'EST PAS, et ou seize colonnes s'etaient trompees :
%  ni une lucarne de toit (it. « lucernari »), ni un volet (cs, lt, lb),
%  ni un meneau ni un montant de chassis (ar, ca, es, oc), ni une
%  imposte fixe, ni le battant de la grande fenetre (bn), ni le
%  « ranma » de la cloison japonaise. Et « framuga » russe nomme
%  l'imposte entiere, quand le francais nomme le VANTAIL : c'est
%  « fortochka » qu'il fallait, comme l'ukrainien avait deja
%  « kvatyrka ».
%
%  LE GENITIF LITUANIEN PRECEDE SON NOM, et il inversait deux fois
%  l'ordre des renvois au tableau 3 — « varpines (5) smaile (7) » pour
%  « la fleche (7) du clocher (5) », « avilio (21) bitutes (22) » pour
%  « les abeilles (22) de la ruche (21) ». Les deux se sont refaites
%  comme le finnois et l'estonien, qui anteposent aussi : relative pour
%  la premiere, ablatif pour la seconde. Le piege n'etait pas dans le
%  lituanien, il etait dans la phrase de depart.
% ===================================================================

%%K t01-apar-0 apar
\VUsaut{2.0mm}
\VUcentre{12.6pt}{120}{AIŠKINAMOJI KNYGELĖ}
\VUsaut{3.2mm}
\VUcentre{8.4pt}{160}{PRIE}
\VUsaut{2.6mm}
\VUtitre{15.6pt}{0}{Delmo pagalbinių lentelių}
\VUsaut{3.4mm}
\VUornamento{9pt}
\VUsaut{5.0mm}
\VUcentre{13.2pt}{90}{PIRMOJI SERIJA}
\VUsaut{3.0mm}
\VUfilet{16mm}
\VUsaut{5.6mm}

%%K t01-apar-1 apar
\VUtitre{15.0pt}{60}{LENTELĖ Nr. 1}
\VUsaut{1.4mm}
\VUtitre{11.4pt}{0}{Mokykla. --- Licėjus. --- Klasė.}
\VUsaut{2.2mm}
\VUfilet{20mm}
\VUsaut{6.4mm}

%%K t01-tit-1 sub
\VUpk{10.2pt}{40}{Bendrasis aprašymas.}
\VUsaut{3.0mm}

%%K t01-01-1 p
1. --- Ši lentelė vaizduoja \VUgras{klasę} \VUgras{licėjuje}. Šioje klasėje yra
aštuoni \VUgras{stalai} \textsuperscript{(18)} ir aštuoni \VUgras{suolai}
\textsuperscript{(32)}. Kiekviename suole sėdi trys mokiniai. \VUgras{Mokytojas}
\textsuperscript{(7)} sėdi ant \VUgras{kėdės} \textsuperscript{(8)} savo
\VUgras{katedroje} \textsuperscript{(6)}.

%%K t01-02-1 p
2. --- Kairėje nuo katedros stovi \VUgras{stiklinė spinta} \textsuperscript{(15)}.
Dešinėje yra \VUgras{juoda lenta} \textsuperscript{(1)} su savo \VUgras{šluoste}
\textsuperscript{(2)} ir \VUgras{kreida} \textsuperscript{(3)}.

%%K t01-02-2 p
Virš katedros kabo \VUgras{sieninės lentelės} \textsuperscript{(9, 11, 12)} ir
\VUgras{žemėlapis} \textsuperscript{(10)}.

%%K t01-03-1 p
3. --- Ne visi mokiniai sėdi savo \VUgras{vietoje} \textsuperscript{(17)}. Jonas
\textsuperscript{(4)} stovi prie lentos; jis laiko \VUgras{tą} šluostę ir trina
\VUgras{geometrines figūras.}

%%K t01-03-2 p
Petras yra prie katedros; jis renka \VUgras{rašto darbus} \textsuperscript{(5)} ir
neša juos mokytojui.

%%K t01-04-1 p
4. --- \VUgras{Šis} mokytojas yra \VUgras{gyvųjų kalbų} mokytojas. Jis moko mokinius
kalbėti, skaityti ir rašyti svetimomis kalbomis \VUgras{(vokiškai, angliškai,
ispaniškai, itališkai, rusiškai} arba \VUgras{prancūziškai)}.

%%K t01-05-1 p
5. --- Klasė labai erdvi ir labai šviesi. Gale yra platus \VUgras{langas} su dviem
\VUgras{orlaidėmis} \textsuperscript{(78)} ir \VUgras{stiklinės durys,} kad ją būtų
galima vėdinti ir apšviesti.

%%K t01-05-2 p
Šioje klasėje nešalta. Viduryje, kad ją šildytų, stovi labai gera \VUgras{krosnis}
\textsuperscript{(46)} su savo \VUgras{vamzdžiu} \textsuperscript{(45)}.

%%K t01-05-3 p
Prie sienos pritvirtintos \VUgras{pakabos} \textsuperscript{(58)}; ant jų kabo
moksleivių kepurės ir paltai.

%%K t01-tit-2 sub
\VUsaut{5.2mm}
\VUpk{10.2pt}{40}{Žaidimų kiemas.}
\VUsaut{3.6mm}

%%K t01-06-1 p
6. --- \VUgras{Laikrodis} \textsuperscript{(63)} rodo devynias valandas ir penkias
minutes.

%%K t01-06-2 p
\VUgras{Pertrauka} \textsuperscript{(76)} ką tik pasibaigė ir pamoka prasideda iš
naujo. Pertraukos vieta yra \VUgras{kieme} \textsuperscript{(75)}. Matome kelis
mokinius, kurie žaidžia. \VUgras{Prižiūrėtojas} \textsuperscript{(74)} juos stebi.

%%K t01-07-1 p
7. --- Kieme dar matome įvairius asmenis, būtent \VUgras{inspektorių}
\textsuperscript{(73)}, mokyklos vedėją (\VUgras{licėjaus direktorių})
\textsuperscript{(71)} ir \VUgras{cenzorių} \textsuperscript{(72)}.

%%K t01-07-2 p
Inspektorius tikrina mokyklas, licėjaus direktorius vadovauja licėjui, cenzorius
prižiūri mokslus.

%%K t01-07-3 p
Ketvirtasis asmuo yra \VUgras{durininkas} \textsuperscript{(77)}. Jis neša po
pažastimi žurnalą, kad įrašytų neatvykusius.

%%K t01-08-1 p
8. --- Kiemo gale yra \VUgras{gimnastikos įrenginiai} \textsuperscript{(70)} ir
dirbtuvė \VUgras{rankų darbams} \textsuperscript{(69)}.

%%K t01-08-2 p
Dešinėje nuo lentelės pastebime \VUgras{muzikos} ir \VUgras{piešimo}
\VUgras{klases}.

%%K t01-noto-1 noto
\VUnotes{143.2mm}{%
(*) \textit{Krikšto vardus} taip pat patariama perrašyti lotyniška forma. Tai
vienintelis būdas išvengti nesuskaičiuojamų skirtybių. Beje, kaip pasirinkti tarp
\textit{Giovanni, Jean, Johann, Jan, Hans, John, Ivan} ir t. t.? Ar kursime atskirą
krikšto vardų žodyną? Imkime iš lotynų kalbos jų vardininką ir sakykime:
\VUgras{Ioannes, Ioanna; Iulius, Iulia; Iakobus; Andreas; Lukas.} (Išsamioji
gramatika).%
}

%%K t01-tit-3 sub
\VUpk{10.2pt}{40}{Išsamusis aprašymas.}
\VUsaut{2.4mm}

%%K t01-tit-4 sub
\VUpk{10.2pt}{40}{Gerieji mokiniai.}
\VUsaut{3.0mm}

%%K t01-09-1 p
9. --- Mokiniai pirmame suole dešinėje nuo katedros yra \VUgras{Henrikas} (*),
\VUgras{Steponas} ir \VUgras{Karolis.}

%%K t01-09-2 p
Henrikas labai dėmesingas ir darbštus; jis klauso mokytojo. Ant savo stalo jis turi
\VUgras{sąsiuvinį} \textsuperscript{(28)}, \VUgras{knygą} \textsuperscript{(29)},
\VUgras{žodyną} \textsuperscript{(30)} ir \VUgras{penalą} \textsuperscript{(31)}. Jo
knyga turi daug \VUgras{puslapių}; kiekviename skyriuje yra keletas
\VUgras{pastraipų}, o kiekvienoje \VUgras{eilutėje} daug \VUgras{žodžių.}

%%K t01-09-4 p
Jo kaimynas Steponas \textsuperscript{(24)} yra uolus. Jis užsirašo į sąsiuvinį kitos
pamokos užduotį. Jis turi gražią \VUgras{rašyseną.} Jo knyga uždaryta. Matome tik
\VUgras{viršelį} \textsuperscript{(27)}. Šalia jo guli jo \VUgras{skriestuvas}
\textsuperscript{(26)}.

%%K t01-09-5 p
Karolis \textsuperscript{(23)} laiko ranka savo knygą; jis ją atverčia, kad dar kartą
perskaitytų savo pamoką. Jis turi, kaip ir kiekvienas mokinys, prieš save
\VUgras{rašalinę} \textsuperscript{(25)}. Jis klusnus.

%%K t01-10-1 p
10. --- Prie gretimo stalo sėdi \VUgras{Liudvikas, Antanas} ir \VUgras{Paulius.}
Liudvikas atsakinėja savo \VUgras{pamoką} \textsuperscript{(22)} ir atsako į mokytojo
\VUgras{klausimus.} Antanas \textsuperscript{(21)} labai nedėmesingas, kartais
mokytojas jį net nubaudžia. Paulius \textsuperscript{(20)} yra geras \VUgras{draugas,}
malonus ir paslaugus. Jis labai dėmesingas.

%%K t01-tit-5 sub
\VUsaut{4.2mm}
\VUpk{10.2pt}{40}{Blogieji mokiniai.}
\VUsaut{3.2mm}

%%K t01-11-1 p
11. --- Už Henriko sėdi \VUgras{Jurgis} \textsuperscript{(38)}. Jis nėra blogas
vaikinas, bet yra \VUgras{plepus,} nedėmesingas ir neramus. Jo
\VUgras{lengvabūdiškumas} ir \VUgras{nepaklusnumas} kelia \VUgras{sumaištį} per
pamoką, ir dažnai jis nusipelno griežto \VUgras{įspėjimo.} Jo \VUgras{juodraštinis
sąsiuvinis} \textsuperscript{(35)} nešvarus, jis \VUgras{aptaško jį rašalu} savo
\VUgras{plunksna} \textsuperscript{(34)}, todėl nuolat naudoja savo \VUgras{trintuką}
\textsuperscript{(36)}; jo \VUgras{portfelis} \textsuperscript{(33)} suplyšęs, nes jis
labai aplaidus. Kodėl jis nešiojasi savo \VUgras{pieštuko laikiklį}
\textsuperscript{(37)} į \VUgras{gyvųjų} kalbų klasę?

%%K t01-11-2 p
Jo kaimynas Edmundas \textsuperscript{(39)} jo nemėgsta. Jis tiesa kiek tingus, bet
tylus ir ramus. Jis neplepa; užuot klausęs, verčiau skaito \VUgras{apsakymus} iš savo
\VUgras{skaitinių knygos.}

%%K t01-11-3 p
Trečiasis šio suolo mokinys yra \VUgras{Liudvikas} \textsuperscript{(40)}. Jis
stropus. Rašo ant balto \VUgras{popieriaus lapo.} Prieš save jis pasidėjo savo
\VUgras{sugeriamąjį popierių} \textsuperscript{(41)}.

%%K t01-12-1 p
12. --- Antrojo suolo mokiniai, kairėje nuo mokytojo, labai jauni. Pirmasis, mažasis
\VUgras{Emilis,} dar nemoka labai gerai rašyti; jis nešiojasi savo \VUgras{grifelinę
lentelę} \textsuperscript{(42)}, savo \VUgras{grifelį} ir savo \VUgras{kempinėlę}
\textsuperscript{(43)}.

%%K t01-12-2 p
Šalia jo yra \VUgras{Augustas} ir \VUgras{Bernardas} \textsuperscript{(44)}. Šis
pastarasis gerai dirba, bet jo \VUgras{apranga} labai apleista.

%%K t01-13-1 p
13. --- Už jų sėdi trys \VUgras{tinginiai. Albertas} nesimoko savo pamokų.
\VUgras{Jokūbas} \textsuperscript{(47)} miega užuot klausęs. Jo \VUgras{tingumas}
užtraukia jam \VUgras{priekaištų} ir net \VUgras{bausmių.} Galiausiai
\VUgras{Kristijonas} \textsuperscript{(48)} per daug nerimtas. Jis blogai
\VUgras{elgiasi} ir nė kiek nesistengia pasitaisyti.

%%K t01-14-1 p
14. --- Jo kaimynai, priešingai, elgiasi gerai. \VUgras{Ernestas}
\textsuperscript{(51)} mėgsta tvarką ir taisyklingumą; jis nuolat naudoja savo
\VUgras{liniuotę} \textsuperscript{(50)} ir savo \VUgras{pieštuką}
\textsuperscript{(49)}. Jis yra \VUgras{eksternas.}

%%K t01-14-2 p
\VUgras{Pranciškus} \textsuperscript{(52)} yra \VUgras{internas}; jis vilki uniformą,
valgo ir miega licėjuje. Jis geras \VUgras{draugas.}

%%K t01-14-3 p
Mauricijus drožia savo \VUgras{pieštuką} savo \VUgras{peiliuku}
\textsuperscript{(56)}; jis labai rūpestingas.

%%K t01-14-4 p
Jis nešiojasi visas savo knygas, savo \VUgras{gramatiką} \textsuperscript{(55)}, savo
\VUgras{užrašų knygelę} \textsuperscript{(54)} ir net \VUgras{rašymo pagalvėlę}
\textsuperscript{(53)}. Jo \VUgras{mokyklinis krepšys} \textsuperscript{(57)} guli ant
grindų už jo.

%%K t01-14-5 p
Du stalus klasės gale užima \VUgras{gerieji mokiniai} \textsuperscript{(19)}.

%%K t01-tit-6 sub
\VUsaut{4.6mm}
\VUpk{10.2pt}{40}{Piešimas ir muzika.}
\VUsaut{3.4mm}

%%K t01-15-1 p
15. --- \VUgras{Piešimo klasėje} ant sienos kabo dailūs \VUgras{pavyzdžiai}, o nuo
lubų nuleista \VUgras{lempa} \textsuperscript{(62)} su \VUgras{gaubtu}.
\VUgras{Mokytojas} \textsuperscript{(60)} labai kantrus; jis taiso mokinių
\VUgras{piešinius} \textsuperscript{(61)} ir moko juos piešti \VUgras{biustą}
\textsuperscript{(59)} iš natūros.

%%K t01-16-1 p
16. --- \VUgras{Muzikos klasė} atrodo labai įdomi. Joje mokomasi skaityti
\VUgras{natas,} solfedžiuoti \VUgras{gamos garsus} \textsuperscript{(67)}, užrašytus
\VUgras{penklinėje} (do, re, mi, fa, sol, la,
si\,=\,C. D. E. F. G. A. B.), pažinti \VUgras{raktus} ir dainuoti žavias
\VUgras{melodijas.} Natos dedamos ant \VUgras{pulto} \textsuperscript{(65)}. Vienas iš
mokinių \VUgras{skambina pianinu} \textsuperscript{(64)}, o \VUgras{muzikos mokytojas}
\textsuperscript{(66)} \VUgras{muša taktą} \textsuperscript{(68)}.

%%K t01-17-1 p
17. --- Pastebime \VUgras{rankų darbų} \textsuperscript{(69)} \VUgras{dirbtuvės}
kampą, kur berniukai gali mokytis obliuoti medį ir dildyti geležį. To yra ne
kiekviename licėjuje.

%%K t01-tit-7 sub
\VUsaut{5.0mm}
\VUpk{10.2pt}{40}{Gamtos mokslų klasė.}
\VUsaut{3.6mm}

%%K t01-18-1 p
18. --- Matematikos mokytojas moko \VUgras{mokinius} \VUgras{geometrijos,
aritmetikos, skaičiavimo} ir \VUgras{metrinės sistemos}. Lentelėje matome pagrindines
geometrines figūras: \VUgras{apskritimą} \textsuperscript{(a)}, \VUgras{kvadratą}
\textsuperscript{(b)}, \VUgras{trikampį} \textsuperscript{(c)}, \VUgras{tiesę}
\textsuperscript{(d)}.

%%K t01-18-2 p
Taip pat matome vieną \VUgras{veiksmą;} tai nėra nei \VUgras{sudėtis,} nei
\VUgras{atimtis,} nei \VUgras{daugyba}: tai iš tiesų \VUgras{dalyba}
\textsuperscript{(e)}.

%%K t01-19-1 p
19. --- Metrinės sistemos lentelėje matome: \VUgras{metrą} \textsuperscript{(a)},
\VUgras{svarstykles} \textsuperscript{(b)} ir jų \VUgras{svarsčius,} \VUgras{litrą}
\textsuperscript{(e)}, \VUgras{dekalitrą} \textsuperscript{(f)}, \VUgras{decilitrą} ir
\VUgras{kubinį metrą} \textsuperscript{(d)}.

%%K t01-19-2 p
Mokiniai taip pat mokosi \VUgras{gamtos mokslų} \textsuperscript{(11)}, zoologijos
\textsuperscript{(a)}, botanikos \textsuperscript{(b)}, geologijos
\textsuperscript{(c)} ir \VUgras{istorijos} \textsuperscript{(12)}.

%%K t01-19-3 p
Spintoje yra \VUgras{fizikos} \textsuperscript{(15)} ir \VUgras{chemijos}
\textsuperscript{(16)} prietaisai.

%%K t01-19-4 p
Ant spintos yra: \VUgras{rutulys} \textsuperscript{(14)}, \VUgras{kūgis} ir
\VUgras{Žemės rutulys} \textsuperscript{(13)}.

%%K t01-19-5 p
Virš lentelių kabo \VUgras{žemėlapis,} vaizduojantis penkias pasaulio dalis:
\VUgras{Europą} \textsuperscript{(g)}, \VUgras{Aziją} \textsuperscript{(e)},
\VUgras{Afriką} \textsuperscript{(f)}, \VUgras{Ameriką} \textsuperscript{(ab)} ir
\VUgras{Okeaniją} \textsuperscript{(h)}, ir du didžiuosius vandenynus ---
\VUgras{Ramųjį} \textsuperscript{(c)} ir \VUgras{Atlanto} \textsuperscript{(d)}.
\VUsaut{6.0mm}
\VUfilet{22mm}
